% !TEX root =./ECE448_Practice_main.tex
%
% L20 practice set (QAM and OFDM). Every question reads a display or works with a
% block parameter: bits per symbol from a constellation, what noise does to
% QPSK, 16-QAM, and 64-QAM, an OFDM spectrum and the link's guard and echo
% sliders, recognizing signals on a waterfall, and choosing a scheme for a use case.
% Blank copy: \begin{solution}[h] reserves h of work space.
% SOLUTIONS copy (\iskey=1) prints the red worked solutions.

\fullwidth{\textbf{Learning Objective 7.1: } I can describe QAM and OFDM and identify
use cases for each.
\\[4pt]\normalfont\small Show your work. Write a \textbf{Documentation} line at the
end (\textbf{None} if you did not collaborate).}

\begin{questions}

%%%%%%%%%%%%%%%%%%%%%%%%%%%%%%%%%%%%%%%%%
\question \textbf{Reading constellations.} A link sends 10,000 symbols per second with RRC
pulses of roll-off $\alpha = 0.25$. Its constellation display shows, at different times,
a grid of 4 by 4 points and a grid of 8 by 8 points, both scaled to the same average power.
\begin{parts}

	\begin{minipage}[t]{0.58\linewidth}
		\part Name each constellation and find its bits per symbol and its bit rate.
		\begin{solution}[1.2in]
		4 by 4 is 16-QAM: $\log_2 16 = 4$ bits per symbol, so 40 kb/s.
		8 by 8 is 64-QAM: $\log_2 64 = 6$ bits per symbol, so 60 kb/s.
		\end{solution}
	\end{minipage}\hfill%
	\begin{minipage}[t]{0.37\linewidth}
		\raggedleft \vspace{12pt}
		16-QAM: \ansbox[1.0in]{40 kb/s}

		\vspace{4pt}
		64-QAM: \ansbox[1.0in]{60 kb/s}
	\end{minipage}

	\begin{minipage}[t]{0.58\linewidth}
		\part Find the occupied bandwidth with each constellation.
		\begin{solution}[1.0in]
		\eq{ B &= (1 + \alpha) R_s = 1.25 \times 10\ \text{kHz} = 12.5\ \text{kHz} }
		The same for both: the bandwidth depends on the symbol rate, not on the number
		of points.
		\end{solution}
	\end{minipage}\hfill%
	\begin{minipage}[t]{0.37\linewidth}
		\raggedleft \vspace{12pt}
		\ansbox[1.0in]{12.5 kHz}
	\end{minipage}

	\part In words, explain how a QAM constellation differs from Lesson 15's 8PSK, and
	why a QAM receiver needs an AGC that a PSK receiver can do without.
		\begin{solution}[1.2in]
		8PSK puts every point on one circle, so only the phase carries data. QAM puts the
		points on a grid, so both the amplitude and the phase carry data. The decision
		between an inner and an outer QAM point depends on the amplitude, so the receiver
		has to hold the signal at the level the constellation expects; a PSK receiver
		decides by phase alone and can ignore the level.
		\end{solution}

\end{parts}

\newpage
%%%%%%%%%%%%%%%%%%%%%%%%%%%%%%%%%%%%%%%%%
\question \textbf{Watching the noise.} You run \texttt{L20-qam-noise.grc} and raise the
Noise amplitude slider, which adds the same noise to the QPSK, 16-QAM, and 64-QAM rows.
\begin{parts}

	\part At one setting, the QPSK display shows four round clouds well apart, the 16-QAM
	display shows sixteen clouds with gaps between them and a few stray points, and the 64-QAM display shows one square
	of scattered dots with no gaps. Describe what a receiver decoding each signal would
	deliver at this setting.
		\begin{solution}[1.3in]
		QPSK: every symbol lands in its own point's region, so the data arrives intact.
		16-QAM: most symbols land in the right region, but the clouds' edges cross into
		neighboring regions now and then, so occasional bits are wrong. 64-QAM: the clouds
		overlap completely, so the receiver picks the wrong point constantly and the data
		is unusable.
		\end{solution}

	\part The noise is the same in all three rows. In words, explain why the three
	displays look so different.
		\begin{solution}[1.2in]
		The noise spreads every point into a cloud of the same size. All three
		constellations have the same average power, so the more points a constellation
		has, the closer together they sit: about 1.41 apart for QPSK, 0.63 for 16-QAM, and
		0.31 for 64-QAM. The same clouds overlap first where the points are closest.
		\end{solution}

	\part A link carries 64-QAM cleanly at short range. In words, describe what you would
	see on its constellation display as the receiver moves away, and what Wi-Fi or LTE
	does at that point.
		\begin{solution}[1.2in]
		The signal weakens relative to the noise, so the 64 clouds grow until they touch
		and merge. Wi-Fi and LTE then step down to a constellation with fewer points,
		16-QAM and then QPSK, and add more error-correcting code, giving up bit rate to
		keep the link working.
		\end{solution}

\end{parts}

\newpage
%%%%%%%%%%%%%%%%%%%%%%%%%%%%%%%%%%%%%%%%%
\question \textbf{An OFDM spectrum and its guard.} A 20 MHz 802.11g Wi-Fi channel uses the
subcarrier layout of \texttt{L20-ofdm-link.grc}: a 64-point FFT, subcarriers $-26$ to $26$
with the one at 0 Hz left empty, and a 16-sample guard. It runs at 20 million samples per
second.
\begin{parts}

	\begin{minipage}[t]{0.58\linewidth}
		\part Find the subcarrier spacing and the width of the occupied band, from the
		center of subcarrier $-26$ to the center of subcarrier $26$.
		\begin{solution}[1.2in]
		\eq{ \Delta f &= 20\ \text{MHz} / 64 = 312.5\ \text{kHz} \\
		52 \times 312.5\ \text{kHz} &= 16.25\ \text{MHz} }
		\end{solution}
	\end{minipage}\hfill%
	\begin{minipage}[t]{0.37\linewidth}
		\raggedleft \vspace{12pt}
		\ansbox[1.0in]{312.5 kHz}

		\vspace{4pt}
		\ansbox[1.0in]{16.25 MHz}
	\end{minipage}

	\begin{minipage}[t]{0.58\linewidth}
		\part Find the length of the guard in microseconds.
		\begin{solution}[0.9in]
		\eq{ 16 / (20 \times 10^6\ \text{samples/s}) = 0.8\ \mu\text{s} }
		\end{solution}
	\end{minipage}\hfill%
	\begin{minipage}[t]{0.37\linewidth}
		\raggedleft \vspace{12pt}
		\ansbox[1.0in]{0.8 $\mu$s}
	\end{minipage}

	\begin{minipage}[t]{0.58\linewidth}
		\part \texttt{L20-ofdm-link.grc} uses the same 16-sample guard at 100,000 samples
		per second. Find the guard's length there, and the delays of its Echo delay slider at
		8 and at 20 samples. Which echo ends inside the guard?
		\begin{solution}[1.2in]
		\eq{ 16 / 100{,}000 &= 160\ \mu\text{s} }
		8 samples is 80 $\mu$s, inside the guard; 20 samples is 200 $\mu$s, past it.
		\end{solution}
	\end{minipage}\hfill%
	\begin{minipage}[t]{0.37\linewidth}
		\raggedleft \vspace{12pt}
		\ansbox[1.3in]{160 $\mu$s}

		\vspace{4pt}
		\ansbox[1.3in]{8 samples (80 $\mu$s)}
	\end{minipage}

	\part In words, describe what \texttt{L20-ofdm-link.grc} showed for an echo inside the
	guard and for one past it: the spectrum, the points after the equalizer, and the text.
		\begin{solution}[1.3in]
		Inside the guard (8 samples): the spectrum rippled, but the equalizer turned the
		points back into a clean 16-QAM grid and every line arrived. Past the guard
		(20 samples): the spectrum rippled more finely, the points after the equalizer
		scattered with no clear grid, and no line arrived intact, although the noise had
		not changed.
		\end{solution}

\end{parts}

\newpage
%%%%%%%%%%%%%%%%%%%%%%%%%%%%%%%%%%%%%%%%%
\question \textbf{Recognizing signals on a waterfall.} For each description, name the
system and whether it uses OFDM or a single carrier, and give one feature that identifies it.
\begin{parts}

	\part A PlutoSDR at 20 MS/s tuned to 2.437 GHz shows short bursts that come and go. Most
	have a flat top about 17 MHz wide with steep edges and a small dip at the center.
		\begin{solution}[1.0in]
		Wi-Fi, OFDM. The bursts are packets; the flat top and steep edges are the
		subcarriers all at the same power, and the dip is the empty subcarrier at 0 Hz.
		\end{solution}

	\part The same PlutoSDR tuned to 740 MHz shows a flat band about 9 MHz wide with steep
	edges, which stays on.
		\begin{solution}[1.0in]
		An LTE downlink, OFDM. A 10 MHz LTE carrier fills about 9 MHz with subcarriers, and
		the base station transmits continuously, unlike Wi-Fi's bursts.
		\end{solution}

	\part A PlutoSDR at 2 MS/s tuned to an FM station shows a rounded signal whose width changes
	with the audio, with a flat shelf of constant width on each side from about 130 to
	200 kHz from its center.
		\begin{solution}[1.0in]
		A hybrid HD Radio station, from Lesson 14: the rounded signal is the analog FM
		host, and the shelves are the OFDM sidebands carrying QPSK.
		\end{solution}

	\part A cable modem uses 256-QAM. In words, explain why you cannot find its signal
	with an SDR and an antenna, and why it can use a single carrier.
		\begin{solution}[1.0in]
		It travels on coaxial cable, which keeps it off the air. The cable also has few
		echoes, so a single carrier with a simple equalizer works and OFDM's guard is not
		needed for that reason.
		\end{solution}

\end{parts}

\newpage
%%%%%%%%%%%%%%%%%%%%%%%%%%%%%%%%%%%%%%%%%
\question \textbf{Choosing a scheme.} For each use case, choose OFDM or a single carrier,
and choose QPSK, 16-QAM, or 64-QAM. Explain each choice in a sentence or two.
\begin{parts}

	\part A 20 MHz-wide data link inside a large metal hangar, where reflections arrive up
	to 1.5 $\mu$s after the direct signal and the signal is strong.
		\begin{solution}[1.1in]
		OFDM with a guard longer than 1.5 $\mu$s. A single carrier filling 20 MHz would send
		symbols well under a tenth of a microsecond long, so each echo would smear across
		dozens of them; with OFDM every
		reflection stays inside the guard and the equalizer removes it. The signal is strong, so the clouds stay small and
		64-QAM can carry the most bits.
		\end{solution}

	\part A telemetry downlink from a small drone at the edge of its range, in open
	country with no reflectors nearby, where the signal is weak.
		\begin{solution}[1.1in]
		A single carrier is enough, because there are few echoes. The signal is weak, so
		QPSK, whose points sit farthest apart, keeps the link working where 16-QAM's
		clouds would merge.
		\end{solution}

	\part A point-to-point link between two towers with dish antennas aimed at each other,
	a clear line of sight, and a strong signal.
		\begin{solution}[1.1in]
		A single carrier, because the narrow beams leave almost no echoes. The strong
		signal supports 64-QAM, as microwave links between towers use.
		\end{solution}

	\part In words, state what decides OFDM or a single carrier, and what decides the
	number of constellation points.
		\begin{solution}[1.0in]
		The echoes decide the first: echoes call for OFDM, with a guard longer than the
		longest echo, while a clean cable or a tower-to-tower link can use one carrier. The noise decides the second: a clean
		signal supports more points and more bits per symbol, and a noisy one needs fewer.
		\end{solution}

\end{parts}

\end{questions}
